\section{Design}

The solution uses one \textbf{goroutine} per \textbf{match} and one \textbf{coordinator} per \textbf{round}.
Each match receives exactly two players and one \textbf{parity function}, then sends one result message back to the round coordinator.
The coordinator collects all results, preserves match order, and builds the \textbf{winner list} for the next round.
The championship coordinator repeats this process until only one player remains.

The implementation is split into three layers:
\begin{itemize}
    \item \textbf{CLI layer}: validates the player count, constructs the initial player slice, and renders the final result;
    \item \textbf{Championship layer}: drives the full elimination and owns the \textbf{round progression};
    \item \textbf{Round layer}: runs \textbf{match workers} concurrently and rebuilds the next player list.
\end{itemize}

The data model stores the tournament outcome explicitly.
Each \textbf{match} stores its round number, match number, the two participants, the winner, and the resulting parity.
Each \textbf{round} stores the ordered list of match results together with the promoted winners.
The \textbf{championship result} stores the final champion and the full sequence of rounds.
Rendering and tests read the stored tournament history instead of recomputing it from the final champion alone.

The coordinator preserves the result order using \texttt{matchIndex}, so completion order and presentation order are
decoupled. This is important because \textbf{goroutines} can finish in any order, but the bracket must remain
deterministic.
